\documentclass[11pt]{article}

\usepackage[T1]{fontenc}
\usepackage[utf8]{inputenc}
\usepackage[english]{babel}
\usepackage[margin=1in]{geometry}
\usepackage{lmodern}
\usepackage{microtype}
\usepackage{mathtools}
\usepackage{amsmath,amssymb,amsthm}
\usepackage{array}
\usepackage{booktabs}
\usepackage{tabularx}
\usepackage{float}
\usepackage{graphicx}
\usepackage{caption}
\usepackage{enumitem}
\usepackage{fancyhdr}
\usepackage[numbers,sort&compress]{natbib}
\usepackage{doi}
\usepackage{xurl}
\usepackage{hyperref}
\usepackage{cleveref}
\usepackage{orcidlink}
\usepackage{titling}



\DeclareMathOperator{\Capacity}{Cap}
\DeclareMathOperator{\Cost}{Cost}
\DeclareMathOperator{\Ran}{Ran}
\DeclareMathOperator{\Ker}{Ker}
\DeclareMathOperator{\tr}{tr}
\DeclareMathOperator{\dist}{dist}

\newcommand{\K}{\mathsf K}
\newcommand{\XiRes}{\Xi}
\newcommand{\Fix}{\operatorname{Fix}}
\newcommand{\preceqq}{\preceq}
 
\newtheorem{theorem}{Theorem}[section]
\newtheorem{lemma}[theorem]{Lemma}
\newtheorem{proposition}[theorem]{Proposition}
\newtheorem{corollary}[theorem]{Corollary}
\theoremstyle{definition}
\newtheorem{definition}[theorem]{Definition}
\newtheorem{remark}[theorem]{Remark}
\floatstyle{ruled}
\newfloat{algorithm}{tbp}{loa}
\floatname{algorithm}{Algorithm}

\urlstyle{same}
\captionsetup{
  font=small,
  labelfont=bf,
  labelsep=period
}
\captionsetup[table]{skip=6pt,position=top}
\captionsetup[figure]{skip=6pt}
\captionsetup[algorithm]{skip=6pt,position=top}
\crefname{algorithm}{algorithm}{algorithms}
\Crefname{algorithm}{Algorithm}{Algorithms}
\crefname{theorem}{theorem}{theorems}
\Crefname{theorem}{Theorem}{Theorems}
\crefname{lemma}{lemma}{lemmas}
\Crefname{lemma}{Lemma}{Lemmas}
\crefname{proposition}{proposition}{propositions}
\Crefname{proposition}{Proposition}{Propositions}
\crefname{corollary}{corollary}{corollaries}
\Crefname{corollary}{Corollary}{Corollaries}
\crefname{definition}{definition}{definitions}
\Crefname{definition}{Definition}{Definitions}
\crefname{remark}{remark}{remarks}
\Crefname{remark}{Remark}{Remarks}
\setlist[enumerate]{leftmargin=*, itemsep=2pt, topsep=4pt}
\setlength{\droptitle}{-3.2em}
\pretitle{\begin{center}\LARGE\bfseries}
\posttitle{\par\end{center}\vspace{0.5em}}
\preauthor{\begin{center}\normalsize}
\postauthor{\par\end{center}\vspace{-0.4em}}
\predate{\begin{center}\small}
\postdate{\par\end{center}\vspace{0.6em}}
\setlength{\emergencystretch}{2em}
\clubpenalty=10000
\widowpenalty=10000
\displaywidowpenalty=10000
\interfootnotelinepenalty=10000
\allowdisplaybreaks[2]
\fancypagestyle{firstpage}{\fancyhf{}
  \fancyfoot[C]{\scriptsize Automorph Inc., Wilmington, DE, USA \quad \texttt{ioannis@automorph.io} \quad \textcopyright\ 2026 \quad CC-BY 4.0}
  \renewcommand{\headrulewidth}{0pt}
  \renewcommand{\footrulewidth}{0pt}
}

\title{Emergence IS the Needle Killer}
\author{Ioannis Tsiokos\,\orcidlink{0009-0009-7659-5964}}
\date{\today}
\hypersetup{
  pdftitle={Emergence IS the Needle Killer},
  pdfauthor={Ioannis Tsiokos},
  hidelinks
}

\begin{document}

\maketitle
\thispagestyle{firstpage}

\begin{abstract}
Many obstruction arguments exhibit a localized failure mode that appears
irreducible at the layer where it is first visible; this paper studies
one audited route by which such needle obstructions cease to be
obstructions, namely passage to a higher layer where additional
admissible dissolving probes are available.  The framework is the
formed-layer membrane setting: zero-energy directions are removed by a
legal quotient, native probes are assigned a quadratic native currency,
unexplained dissolving currency is measured by an adequacy residual, and
predictive membranes, dissolving transport records, and the all-six
taxonomy control the transfer.  The layer-dissolving endpoint theorem
shows that, under predictive native bounds, residual bounds, and
transport hypotheses, the dissolving currency is bounded uniformly
across accepted stages, so no layer-dissolving predictive needles exist
within the declared scope.  The complementary duality-confinement
membrane theorem shows that, under a separating anti-invariant readout
and a vanishing-trace sequence of domination records, visible object
mass is confined to the fixed locus of the involution.  These are
framework-bound dissolution results, not a universal
singularity-elimination theorem or a claim about any named external
problem.  The adequacy-residual calculus on which the layer-dissolving
endpoint depends is developed in \citet{Tsiokos2026XiCompanion}.
\end{abstract}

\medskip
\noindent\textbf{Keywords.} Six Birds theory; formed-layer membrane; emergence; dissolution; adequacy residual; duality-confinement; all-six taxonomy.



\section{Introduction}
\label{sec:intro}

Many obstruction arguments begin in the layer where the obstruction is
first visible.  One has a space of objects, a family of tests, and a
localized failure mode that survives those tests.  A \emph{lower layer}
will mean this original mathematical setting, including the tests and
bookkeeping available there.  A \emph{native probe family} is the
family of linear measurements available in that lower layer.  A
\emph{needle obstruction}, or simply a \emph{needle} after first use, is
a localized obstruction, singular concentration, or persistent failure
mode that cannot be removed using the native probes of the layer in
which it appears.

This terminology is deliberately neutral.  The paper does not need a
single classical model of singularity in order to begin: the same lower
layer may present a point concentration, a support-overlap failure, a
direction outside a measurement range, or a symmetry-breaking
displacement.  The shared shape is what matters: the obstruction is
localized enough to be named, persistent enough to survive the tests
that the layer authorizes, and sharp enough that improving constants
inside the same probe family does not remove it.

The difficulty is that persistence under native probes is not the same
as absolute irreducibility.  A point concentration may survive all
energy-shell tests because those tests average over the very coordinate
in which the concentration lives \citep{Lions1984ConcentrationI,Evans2010PartialDifferential}.  A residual vector may be invisible to
a chosen family of linear measurements, even though it is detected
immediately after one adds a missing readout.  A displacement may look
like genuine motion when only symmetric probes are allowed, although the
anti-invariant part of the same displacement is the quantity that must
be audited.  In each case the native layer has not necessarily found an
unavoidable singularity; it may have found the limit of its own probe
language.

This paper studies one audited route by which that situation is
resolved.  \emph{Emergence} means passage from the original setting to a
richer setting with additional admissible probes and records that can
measure or eliminate the obstruction.  The \emph{higher layer} is that
richer mathematical setting, equipped with the extra probes and
bookkeeping needed for the argument.  An \emph{emergent layer} is a
higher layer viewed as the destination of such a passage.  A
\emph{dissolving probe family} is a family of measurements available in
the higher layer that may see an obstruction the native probes cannot
explain.  \emph{Dissolution} is the conversion of a lower-layer needle
into bounded residual accounting, confinement to a fixed locus, or a
no-critical-pair conclusion after passing to the richer layer.

``Dissolution'' is meant literally in this restricted sense: the
obstruction is not declared harmless by changing notation, nor ignored
by embedding the problem into a larger ambient space.  The higher
layer must carry probes that are admissible for the original energy
and records that explain how the lower-layer statement is being
transported.  A finite-dimensional sketch makes the discipline
concrete.  Suppose a semidefinite energy prices vectors, a native
family records only a few responses, and a vector direction remains
unpriced because it lies outside the native explanation range.  If a
higher response space adds a concentration or anti-invariant readout,
the formerly hidden direction becomes an audited residual, and the
question is no longer whether the native probes alone can remove it,
but whether the residual is bounded by the declared higher-layer
budget.

The proof apparatus is the \emph{formed-layer membrane}: a package of
quotient, probe, currency, residual, and budget data used to test
whether bounds transfer across layers.  The first step is the
\emph{legal quotient}, the quotient or restricted subspace obtained
after removing zero-energy directions that should not count as
admissible degrees of freedom.  On this quotient the native probes carry
a \emph{native currency}, a quadratic accounting object that measures
what the lower-layer probes can spend and explain.  The higher-layer
dissolving probes then introduce an \emph{adequacy residual}, the part
of the dissolving currency not explained by the native currency.  The
residual is the gate at which the present development meets the
adequacy-residual calculus of \citet{Tsiokos2026XiCompanion}, where the
projection, optimality, exactness, transport, and blind-spot witness
theorems are developed.

The ingredients are standard in their local mathematical behavior.
Quadratic forms price vectors, quotienting removes directions that cost
zero energy, positive semidefinite matrices are compared by Loewner
order, and residual blocks are measured by Schur-complement accounting.
The contribution of the framework is not to rename these constructions,
but to keep their roles synchronized across a passage between layers.
The quotient must be the quotient used by the probes; the native
currency must be the one used by the predictive budget; the residual
must be the unexplained part seen by the dissolving probes; and the
endpoint theorem must spend exactly those currencies, not a nearby
informal analogue.

With this accounting in place, the membrane results have a direct
interpretation.  A \emph{predictive native membrane} is a native-layer
budget condition whose failure is equivalent to an accepted
critical-pair witness.  Thus the lower layer is not merely asserted to
be controlled; the accepted predictive family carries finite Loewner
bounds that can be propagated to the dissolution step.  The
\emph{layer-dissolving endpoint} is the theorem that transports native
and residual bounds to a higher-layer dissolving budget.  It is the
first of the paper's two main results: what looked like a needle at the
lower layer is bounded after the dissolving probes and adequacy residual
are included.

The second main result is \emph{duality-confinement}, the membrane
theorem saying that, when the relevant anti-invariant currency is
forced into a vanishing budget, visible objects are confined to the
fixed locus of the involution.  This is not redundant with the
layer-dissolving endpoint.  The endpoint transports bounded dissolving
currency into a target budget; duality-confinement treats the case in
which the moving, anti-invariant part of the object is squeezed away by
the available ledger.  The two mechanisms address different ways in
which a lower-layer obstruction can cease to be an obstruction once the
right higher-layer records are present.

In the formed-layer membrane setting, emergence can dissolve
needle-class lower-layer obstructions when the legal quotient, native
currency, adequacy residual, predictive membrane bounds, and dissolving
transport records are all present.  \Cref{sec:dissolution} proves the
layer-dissolving membrane endpoint theorem, the transport statement that
places the dissolving currency inside the target budget.
\Cref{sec:duality-confinement} proves the duality-confinement membrane
theorem, the fixed-locus conclusion obtained from anti-invariant ledger
collapse.  \Cref{sec:taxonomy} gives the qualifying discipline: the
\emph{all-six taxonomy}, a six-role checklist of rewrite, feasibility,
routing, staging, packaging, and audit data, determines when a
dissolution instance is a valid transfer record rather than an informal
analogy.

That last qualification is essential.  A higher-layer argument can fail
by using the right-looking probes in the wrong quotient, by transporting
data through an unrecorded map, by proving a budget at the wrong stage,
or by hiding a residual defect in the packaging.  The all-six taxonomy
does not classify all possible needles.  It records the roles that must
be present for the instances treated here, so that the conclusion is a
checked transfer from lower-layer obstruction to higher-layer
dissolution.

The scope is correspondingly narrow.  The paper proves a formed-layer
dissolution mechanism for needle-class obstructions under its standing
records; it is not a universal singularity-elimination theorem and does
not claim that every mathematical obstruction can be removed by moving
to a higher layer.  It also does not give an application proof for the
Riemann hypothesis \citep{Bombieri2000Riemann}, Navier--Stokes regularity
\citep{Fefferman2000NavierStokes}, or any other named external problem
in this submission.  Those applications would require
separate problem-specific records and separate verification of the
required membrane hypotheses.

The organization is as follows.  \Cref{sec:framework} constructs the
legal quotient, native probe accounting, and minimum-spend language used
throughout the paper.  \Cref{sec:xi-interface} states the adequacy gate
that separates native explanation from residual dissolving currency.
\Cref{sec:predictive-membranes} develops the predictive native membrane
criterion and the finite propagation bounds that feed the endpoint.
\Cref{sec:dissolution} proves the layer-dissolving endpoint, and
\Cref{sec:duality-confinement} proves the complementary
duality-confinement theorem.  \Cref{sec:critical-pairs} records the
strict-extension completion needed for critical-pair control, while
\Cref{sec:taxonomy} closes the body by stating the all-six and
adequacy-aware transfer discipline.  The discussion and conclusion then
return to the scope of the dissolution claim and the role of the
residual calculus of \citet{Tsiokos2026XiCompanion}.

 

\section{The Formed-Layer Framework}
\label{sec:framework}

The introduction used the language of needles, emergence, and
dissolution informally.  We now pass to the finite-dimensional
operator setting in which those words acquire a budget.  The section
first fixes the legal quotient and packaged native currency, then
turns that currency into a minimum-spend formula, and finally records
the Schur-block tests that certify lawful witnesses.

\subsection{Legal quotient and packaged native currency}

Let \(E\) be a finite-dimensional carrier, meaning the ambient vector
space on which the lower-layer audit is first written.  An
\emph{audit energy} is a positive semidefinite quadratic form
\[
  C=C^*\succeq 0
\]
on \(E\).  Since \(C\) may have a kernel, formulas that look like
division by \(C\) must be interpreted only after the zero-energy
directions have been removed.  The \emph{legal quotient} is represented
through the legal carrier
\[
  E_C=(\Ker C)^\perp,\qquad C_0=C|_{E_C}>0.
\]
Thus \(E_C\) is the positive-energy representative of the quotient by
\(\Ker C\), and \(C_0\) is the strictly positive restriction on which
pseudoinverse and inverse-style capacity formulas have their intended
variational meaning.

A \emph{native probe family} is the family of linear measurements
available in the original layer.  We write it as a linear map
\[
  L:E\to Y,
\]
where \(Y\) is the response space.  The probe family is
\emph{null-legal} when
\[
  L(\Ker C)=0.
\]
Equivalently, the measurements vanish on zero-energy directions, so
they descend to the legal quotient.  We use ``null-legality'' and
``null-mode legality'' interchangeably for this condition.  It is not
cosmetic: if null-legality fails, a matrix expression using the
Moore--Penrose pseudoinverse \(C^\dagger\)
\citep{Penrose1955GeneralizedInverse,BenIsraelGreville2003GeneralizedInverses}
can appear finite because the pseudoinverse suppresses
the kernel, while the true variational capacity is infinite, since a
zero-energy direction would change the response at no energy cost.  The formed-layer
framework therefore treats legality as part of the data, not as a
regularity condition to be checked after the fact.

In applications the audited sub-problem may be expressed through an
internal package rather than through the full carrier \(E\).  An exact
package \(\Gamma\) is the package data selecting the admissible
directions and coordinates used for that audited sub-problem, without
changing the probe and energy information that the sub-problem is meant
to preserve.  It determines the \emph{packaged audit} and
\emph{packaged native family}
\[
  C_\Gamma=\Gamma^*C\Gamma,\qquad L_\Gamma=L\Gamma.
\]
The following definition is the first currency object used throughout
the paper.

\begin{definition}[Packaged native currency]\label{def:main:packaged-currency}
Let \(\Gamma\) be an exact package.  Set
\[
  C_\Gamma=\Gamma^*C\Gamma,\qquad L_\Gamma=L\Gamma.
\]
If \(L_\Gamma\) is null-legal for \(C_\Gamma\), the
\emph{packaged native currency} is
\[
  \K_{C,\Gamma}(L)=L_\Gamma C_\Gamma^\dagger L_\Gamma^*.
\]
\end{definition}

This matrix is the quadratic response currency attached to the native
probe family after the package and legal quotient have been imposed.
It is the object against which later native membrane budgets are
compared.

\begin{theorem}[Probe-family currency]\label{thm:main:probe-family-currency}
Assume the relevant null-mode legality hypotheses for the packaged
native family.  Then, for every native recombination \(y\in Y\),
\[
  \Capacity_C(y^*L;\Ran\Gamma)=y^*\K_{C,\Gamma}(L)y.
\]
Consequently, a native membrane budget is the Loewner inequality
\[
  \K_{C,\Gamma}(L)\preceqq\Theta,
\]
and failure of this budget is equivalent to the existence of a
critical-pair witness \(y\) such that
\[
  y^*(\K_{C,\Gamma}(L)-\Theta)y>0.
\]
\end{theorem}

\begin{proof}
The recombination \(y\) turns the native probe family into the scalar
probe \(y^*L\).  Restricting the admissible vectors to \(\Ran\Gamma\)
is the same as evaluating this scalar probe through the packaged data
\((C_\Gamma,L_\Gamma)\).  Null-mode legality ensures that this
evaluation descends to the positive-energy quotient, so the standard
variational capacity calculation gives the quadratic form
\[
  y^*L_\Gamma C_\Gamma^\dagger L_\Gamma^*y
  =y^*\K_{C,\Gamma}(L)y.
\]
In the general null-legal regime, this identity gives \(\K_{C,\Gamma}(L)\)
its capacity interpretation: it is not merely a formal pseudoinverse
expression, but the finite minimum-spend value of the native
recombination.

The budget statement is then the elementary separation property of the
Loewner order \citep{Lowner1934Monotone,Bhatia1997MatrixAnalysis}.  The inequality
\(\K_{C,\Gamma}(L)\preceqq\Theta\) holds exactly when
\[
  y^*(\K_{C,\Gamma}(L)-\Theta)y\le 0
\]
for every response vector \(y\).  If the inequality fails, some \(y\)
has positive excess; that vector is the critical-pair witness for the
native budget failure.\footnote{Mechanized in Lean as
\texttt{SixBirdsNeedles.Main.LegalQuotient.probeFamilyCurrency}; the Lean
theorem proves the point-evaluation matrix identity under stricter
positivity on a fixed positive legal subspace.  The full statement is
proved here.  See \Cref{subsec:paper1-per-label-coverage}.}
\end{proof}

\subsection{Minimum-spend duality}

The preceding currency identity says what a native recombination can
extract from admissible inputs, and it packages native budgets as
Loewner inequalities such as \(\K_{C,\Gamma}(L)\preceqq\Theta\).
The dual question is how much audit energy is needed to realize a
specified native response.  The \emph{minimum spend} for a desired
response \(z\) is the variational infimum of the packaged audit energy
over native inputs whose probe value is \(z\), whenever such inputs
exist.

\begin{theorem}[Minimum-spend duality]\label{thm:main:minimum-spend-duality}
For a desired response \(z\), define
\[
  \Cost_C(z;L,\Gamma)
    :=\inf\{a^*C_\Gamma a: L_\Gamma a=z\}.
\]
Let \(L_{\Gamma,0}\) denote \(L_\Gamma\) restricted to the legal
subspace of \(C_\Gamma\), after quotienting by \(\Ker C_\Gamma\).
Then
\[
  \Cost_C(z;L,\Gamma)
    =z^*\K_{C,\Gamma}(L)^\dagger z
\]
when
\[
  z\in\Ran L_{\Gamma,0}=\Ran\K_{C,\Gamma}(L),
\]
and \(\Cost_C(z;L,\Gamma)=+\infty\) otherwise.
\end{theorem}

In the proof below, \(C_{\Gamma,0}\) denotes the positive restriction
of \(C_\Gamma\) to the legal subspace used by \(L_{\Gamma,0}\).

\begin{proof}
On the legal subspace of \(C_\Gamma\), the packaged audit is positive
and \(L_{\Gamma,0}\) is an ordinary finite-dimensional linear map from
the audited carrier to the response space.  If
\(z\in\Ran L_{\Gamma,0}\), the admissible set in the definition of
\(\Cost_C(z;L,\Gamma)\) is nonempty.  Among all preimages of \(z\), the
minimizer is the least-energy preimage, equivalently the least-norm
preimage after conjugating the legal carrier by the square root of the
positive audit.  The standard pseudoinverse variational identity
\citep{BoydVandenberghe2004Convex} then
identifies the achieved energy with
\[
  z^*(L_{\Gamma,0}C_{\Gamma,0}^{-1}L_{\Gamma,0}^*)^\dagger z
  =z^*\K_{C,\Gamma}(L)^\dagger z.
\]
This is the finite-range branch of the formula.

If \(z\notin\Ran\K_{C,\Gamma}(L)=\Ran L_{\Gamma,0}\), then no legal
native input has packaged probe value \(z\).  Null-legality ensures
that adding a vector in \(\Ker C_\Gamma\) cannot change the response,
so there is no hidden zero-energy correction that would realize
\(z\).  The admissible set in the defining infimum is empty, and the
minimum spend is therefore \(+\infty\).\footnote{Mechanized in Lean as
\texttt{SixBirdsNeedles.Main.LegalQuotient.minimumSpendDuality}; the Lean
theorem proves the off-range infinite-cost branch under the finite
packaged legal-carrier hypotheses.  The full statement is proved here.
See \Cref{subsec:paper1-per-label-coverage}.}
\end{proof}

Minimum-spend duality gives the packaged native currency its cost
interpretation.  The matrix \(\K_{C,\Gamma}(L)\) records the finite
responses that native probes can explain, while its pseudoinverse
prices the least audit energy needed to produce each explained
response; responses outside its range are not expensive but unavailable
to the native layer.

\subsection{Schur squares and lawful witnesses}

Before moving to the interface section, we need the block algebra that
separates a component already represented by one block from the
remaining Schur margin.  The classical Schur-complement test
\citep{Schur1917Potenzreihen,Zhang2005SchurApplications,HornJohnson2013MatrixAnalysis}
has to be read through the legal quotient introduced above.  When the lower block
is only semidefinite, the off-diagonal data must itself be legal;
otherwise the pseudoinverse gives a finite-looking number that is not a
finite variational cost.

The scalar Schur block criterion is the basic case.  Let
\(C=C^*\succeq 0\) be a positive semidefinite quadratic form on \(E\),
let \(a\in\mathbb R\), and let \(g\in E\).  Then
\[
  \begin{pmatrix}
    a & g^*\\
    g & C
  \end{pmatrix}\succeq 0
\]
if and only if
\[
  g\perp\Ker C
  \qquad\text{and}\qquad
  g^*C^\dagger g\le a.
\]
The condition \(g\perp\Ker C\) is the null-legality condition for the
off-diagonal vector.  If it fails, the displayed pseudoinverse
expression still has a formal value, but the block cannot be positive:
a zero-energy direction would see a nonzero off-diagonal coupling.  This
is the scalar/vector instance of the lawful Schur identity used below;
the corresponding positive-semidefinite block criterion is the
classical pseudoinverse condition of \citet{Albert1969PSDPartitioned}.

The same idea appears in compressed form.  An \emph{internal low
package} \(\Phi\) is a linear embedding that selects an interior
subspace of the relevant carrier; it is dual in role to the exact
package \(\Gamma\) above, but it selects an interior block rather than
restricting to an admissible carrier.  Let \(M\succeq 0\) be a positive
semidefinite form on a larger Hilbert space \(H\), and let
\(Q:H\to H\) be the orthogonal projection onto the orthogonal complement
of \(\Ran\Phi\).  Define
\[
  A=\Phi^*M\Phi,\qquad
  B=QM\Phi,\qquad
  C_Q=QMQ|_{QH}.
\]
When \(B\) is null-legal for \(C_Q\), in the sense that it vanishes on
the zero-energy directions of the complementary block, the compressed
Schur margin is
\[
  S=A-B^*C_Q^\dagger B.
\]
Under the same quotient-legality hypotheses, this margin satisfies
\[
  S=(\Phi^*M^\dagger\Phi)^{-1}.
\]
Thus the compressed Schur margin records the part of \(M\) carried by
the interior subspace after the complementary part has been subtracted.
Null-legality is what makes both displayed inverse operations
variationally meaningful: \(C_Q^\dagger\) is applied only after its
zero-energy directions have been quotiented away, and the inverse of
\(\Phi^*M^\dagger\Phi\) is taken only on the resulting positive
interior block.

This is the algebraic setting for a \emph{lawful witness}.  A witness
vector for a Loewner failure such as
\(\K_{C,\Gamma}(L)\not\preceqq\Theta\) is lawful when it satisfies the
null-legality constraints inherited from the legal quotient, so that it
can be packaged into a Schur block of the form above without hiding an
infinite variational cost behind a finite pseudoinverse expression.
Lawful witnesses are the audit-respecting analogue of the witnesses
produced by the scalar Schur criterion.

The legal quotient and null-legality requirements introduced above
are what make the minimum-spend formula meaningful.  The scalar and
compressed Schur criteria give the algebraic certificates by which a
candidate witness is checked against those requirements before the
next section introduces the interface statement.

 

\section{The Adequacy-Residual Interface}
\label{sec:xi-interface}

This paper uses the adequacy residual only as a gate.  It is the Schur
leftover between dissolving and native currencies on the legal quotient,
invoked here to give an upper bound on the dissolving block in terms
of native and residual budgets.  The full residual calculus---projection, optimality,
exactness in multiple equivalent forms, transport, post-processing, and
witness characterization---is developed in
\citet{Tsiokos2026XiCompanion}.  This section records only the interface
identity needed later in the paper.

The point of making this gate explicit is to prevent the higher-layer
probe family from being treated as a free enlargement of the native
data.  A dissolving probe may see more than a native probe sees, but the
extra visibility has to be paid for by a residual budget.  The interface
below separates the part explained by the native block from the part
that must be carried forward as residual currency.

Recall the native probe family \(L:E\to Y\) from
\Cref{sec:framework}.  A \emph{dissolving probe family}, also called a
\emph{layer-dissolving probe family} when it is used in the endpoint
theorem, is a family of measurements available in the higher layer that
may see information the native probes cannot explain; write it as
\(D:E\to Z\).  On the legal carrier \(E_C\), with
\(C_0=C|_{E_C}>0\) and with \(L_0,D_0\) denoting the corresponding
restrictions, set
\[
  K_{LL}=L_0C_0^{-1}L_0^*,\quad
  K_{DL}=D_0C_0^{-1}L_0^*,\quad
  K_{DD}=D_0C_0^{-1}D_0^*.
\]
The block \(K_{LL}\) is the native self-currency, \(K_{DL}\) is the
dissolving-by-native cross currency, and \(K_{DD}\) is the dissolving
self-currency.  We also write \(K_{LD}=K_{DL}^*\) for the adjoint cross
block.

Together these are the blocks of the joint currency of the native and
dissolving probe families after all zero-energy directions have been
removed.  The notation mirrors the Schur-block discussion in
\Cref{sec:framework}: the native block is the part already priced by the
lower layer, the dissolving block is the higher-layer quantity later to
be bounded, and the cross block records how much of the dissolving
response is routed through native responses.

The \emph{adequacy residual} is the Schur leftover
\[
  \XiRes_C(D\mid L)=K_{DD}-K_{DL}K_{LL}^\dagger K_{LD}.
\]
Thus the residual is not an arbitrary error term.  It is the part of the
dissolving self-currency that remains after the native block has
accounted for everything it can through the cross currency, and it is
the quantity that will receive a separate budget here.  The formal
development of the residual as a calculus---projection identity,
optimal-residual identity, exact-adequacy equivalence, chain rule,
post-processing laws, and blind-spot witnesses---is the body of
\citet{Tsiokos2026XiCompanion}; the present paper cites that calculus by
name rather than reproving its theorems.

In the interface estimate below, \(\Theta_Y\) denotes a budget on the
native response space \(Y\), while \(\Xi_Z\) denotes a residual
budget on the dissolving response space \(Z\).

\begin{theorem}[Adequacy interface]\label{thm:main:adequacy-interface}
Let \(A_*:=K_{DL}K_{LL}^\dagger\).  Then
\[
  K_{DD}=A_*K_{LL}A_*^*+\XiRes_C(D\mid L).
\]
Consequently, if \(K_{LL}\preceqq\Theta_Y\) and
\(\XiRes_C(D\mid L)\preceqq\Xi_Z\), then
\[
  K_{DD}\preceqq A_*\Theta_YA_*^*+\Xi_Z.
\]
Moreover,
\[
  \XiRes_C(D\mid L)=0
  \iff \exists A\text{ with }D_0=AL_0
  \iff \Ker L_0\subseteq\Ker D_0.
\]
\end{theorem}

\begin{proof}
The first assertion is the Schur identity prepared by
\Cref{sec:framework}.  By definition of \(A_*\), the native-explained
part of the dissolving block is
\[
  A_*K_{LL}A_*^*
  =K_{DL}K_{LL}^\dagger K_{LD},
\]
where the pseudoinverse is taken on the native currency coming from the
legal quotient.  Hence
\[
  K_{DD}
  =A_*K_{LL}A_*^*
   +\bigl(K_{DD}-K_{DL}K_{LL}^\dagger K_{LD}\bigr)
  =A_*K_{LL}A_*^*+\XiRes_C(D\mid L).
\]
The point of the legal-quotient discipline is that \(K_{LL}^\dagger\)
is not a free formal inverse: it is the pseudoinverse of the native
currency after null-legality has forced the relevant probes to descend
to the quotient.  This is the same legality issue as in
\Cref{sec:framework}: without it, the Schur expression could look finite
while the variational problem it is meant to represent has an illegal
zero-energy direction.

The budget consequence is monotonicity in the Loewner order.  If
\(K_{LL}\preceqq\Theta_Y\), then conjugating by \(A_*\) gives
\[
  A_*K_{LL}A_*^*\preceqq A_*\Theta_YA_*^*.
\]
Adding the residual bound
\(\XiRes_C(D\mid L)\preceqq\Xi_Z\) to this inequality and using the
Schur splitting gives
\[
  K_{DD}\preceqq A_*\Theta_YA_*^*+\Xi_Z.
\]
This is the interface estimate used later by the layer-dissolving
endpoint theorem in \Cref{sec:dissolution}.  Its two inputs are
exactly the native budget on \(K_{LL}\) and the residual budget on
\(\XiRes_C(D\mid L)\).

Finally, \(\XiRes_C(D\mid L)=0\) is the exact-adequacy case.
If \(D_0=AL_0\), then every vector killed by \(L_0\) is killed by
\(D_0\), so \(\Ker L_0\subseteq\Ker D_0\), and the dissolving currency
has no Schur leftover after the native block explains it.  Conversely,
in finite dimension, the inclusion
\(\Ker L_0\subseteq\Ker D_0\) says exactly that \(D_0\) factors through
the quotient by \(\Ker L_0\), hence through the native map \(L_0\);
this gives a linear map \(A\) with \(D_0=AL_0\).  This proves the
three-form exactness statement used here.  The full optimal native
explanation characterization, including the Loewner-minimal residual
form and the additional equivalent formulation, belongs to
\citet{Tsiokos2026XiCompanion}.

The exact case marks the boundary between native sufficiency and
residual currency that still has to be budgeted.  In the nonzero case,
the present paper does not classify the residual or produce witnesses
for it; that classification is developed in
\citet{Tsiokos2026XiCompanion}.  We only need the residual as a positive
budgeted object that can be transported into the dissolution theorem.\footnote{Mechanized in Lean as
\texttt{SixBirdsNeedles.Main.XiInterface.adequacyInterface}; the Lean
theorem proves the Schur identity under explicit algebraic hypotheses
on the legal subspace and packaged data.  The full statement, including
Loewner monotonicity and exact-factorization equivalence, is proved
here.  See \Cref{subsec:paper1-per-label-coverage}.}
\end{proof}

This theorem is the sole interface the rest of the paper needs to the
\(\XiRes\)-machinery.  The predictive native membranes of
\Cref{sec:predictive-membranes} spend against \(\K\) and \(\Theta_Y\);
the layer-dissolving endpoint of \Cref{sec:dissolution} consumes a
bounded \(\XiRes_C(D\mid L)\) and produces the dissolving budget for
\(K_{DD}\); and the duality-confinement theorem of
\Cref{sec:duality-confinement} imports the residual budget without
re-developing it.  Keeping the interface narrow lets the present paper
use the residual calculus without duplicating it.  The blind-spot
interpretation---that the adequacy residual captures the part of the
dissolving probes' measurements that the native probes cannot
explain---is developed quantitatively as blind-spot currency in
\citet{Tsiokos2026XiCompanion}.

 

\section{Predictive Native Membranes}
\label{sec:predictive-membranes}

The adequacy interface of \Cref{sec:xi-interface} asks for two inputs:
a native budget on the native block and a residual budget on the
adequacy residual.  This section supplies the native side in staged
form.  At each accepted stage of a refinement or acceptance process, a
predictive currency records the native-layer measurement that will later
be spent by the layer-dissolving endpoint.

\subsection{Predictive currency and the membrane criterion}

In many applications the native data is not accepted all at once.  It is
assembled stage by stage: by refinement, enlargement, or passage to an
accepted stage.  At an accepted stage \(j\), we have a stage-local
instance of the framework from \Cref{sec:framework}: a carrier, an audit
energy \(C_j\), an exact package \(\Gamma_j\), a native probe family
\(L_j\) with its response space, and a routing map \(R_j\) into a common
response target.  The predictive currency at stage \(j\) is the
packaged native currency for these data, transported by \(R_j\).

Informally, the predictive currency is
\[
  \widehat K_j
  =R_jL_j\Gamma_j(\Gamma_j^*C_j\Gamma_j)^\dagger
    \Gamma_j^*L_j^*R_j^*.
\]
The middle block
\[
  L_j\Gamma_j(\Gamma_j^*C_j\Gamma_j)^\dagger\Gamma_j^*L_j^*
\]
is the stage-local instance of the packaged native currency
\(\K_{C_j,\Gamma_j}(L_j)\) from \Cref{sec:framework}; the outer
conjugation by \(R_j\) places that currency in the common response
space.  As before, this expression is variationally meaningful only
when \(L_j\Gamma_j\) is null-legal for the packaged audit
\(\Gamma_j^*C_j\Gamma_j\), which is the legal-quotient discipline of
\Cref{sec:framework} applied at stage \(j\).

\begin{theorem}[Predictive native membrane]\label{thm:main:predictive-native-membrane}
A native predictive membrane with budget \(\Theta\) is the bound
\[
  \widehat K_j\preceqq\Theta\qquad\forall j\text{ accepted}.
\]
Equivalently, no accepted stage contains a native recombination witness
\((j,y)\) with
\[
  y^*(\widehat K_j-\Theta)y>0.
\]
\end{theorem}

\begin{proof}
The equivalence is the Loewner-separation argument from
\Cref{sec:framework}, now indexed by accepted stages.  For a fixed
stage \(j\), the inequality \(\widehat K_j\preceqq\Theta\) holds if and
only if
\[
  y^*(\widehat K_j-\Theta)y\le 0
\]
for every response vector \(y\).  Requiring the inequality for every
accepted \(j\) is therefore the same as ruling out every pair \((j,y)\)
with strict positive excess; these are precisely the native critical-pair
witnesses for the predictive currency.\footnote{Verified in Lean as
\texttt{SixBirdsNeedles.Main.PredictiveNative.predictiveNativeMembrane};
see \Cref{subsec:paper1-per-label-coverage}.}
\end{proof}

Thus a predictive native membrane is a stage-by-stage Loewner
certificate: producing one is exactly ruling out all accepted
critical-pair witnesses.  The next step is to propagate such certificates
across the accepted refinement.

\subsection{Summable predictive propagation}

The preceding criterion is stagewise: it certifies a Loewner budget at
each accepted stage.  The remaining question is when such certificates
survive passage from one accepted stage to the next, and when they can
be aggregated over the whole accepted refinement in a common response
space.  The mechanism is multiplicative inflation at each step, together
with additive defects whose total contribution is summable.

\begin{theorem}[Summable predictive propagation]\label{thm:main:summable-predictive-propagation}
If, in a common response space,
\[
  \widehat K_{j+1}\preceqq (1+\varepsilon_j)\widehat K_j+D_j,
\]
with
\[
  \prod_j(1+\varepsilon_j)<\infty,\qquad \sum_jD_j\preceqq D_\infty,
\]
then
\[
  \widehat K_n\preceqq P_\infty(\widehat K_0+D_\infty),
\]
where \(P_\infty\) denotes \(\prod_j(1+\varepsilon_j)\).
\end{theorem}

\begin{proof}
The hypothesis is the Loewner inflation-plus-defect step.  Iterating it
from stage \(0\) to stage \(n\) gives the finite-stage envelope
\[
  \widehat K_n
  \preceqq
  \Bigl(\prod_{i=0}^{n-1}(1+\varepsilon_i)\Bigr)\widehat K_0
  +
  \sum_{i=0}^{n-1}
  \Bigl(\prod_{k=i+1}^{n-1}(1+\varepsilon_k)\Bigr)D_i .
\]
Each scalar product is bounded above by \(P_\infty\), and the defect
terms are positive semidefinite.  Hence the finite sum of propagated
defects is bounded by \(P_\infty D_\infty\) in the Loewner order, while
the propagated initial term is bounded by \(P_\infty\widehat K_0\).  This
is exactly the displayed estimate.

The conclusion is uniform in the finite stage index: the constants
\(P_\infty\) and \(D_\infty\) are fixed by the propagation data, so the
same envelope bounds every \(\widehat K_n\) for which the iteration is
defined.  No convergence of the sequence \(\widehat K_n\) is asserted or
needed.

The limiting character of the statement lies in the hypotheses, not in
the conclusion.  The infinite product and the operator-valued sum are
read as uniform bounds over all finite truncations: finite products are
bounded by \(P_\infty\), and finite defect sums are bounded by
\(D_\infty\) in the Loewner order.  This is the passage from the
infinite-index hypotheses to the uniform finite-truncation bounds used
above.\footnote{Mechanized in Lean as
\texttt{SixBirdsNeedles.Main.PredictiveNative.summablePredictivePropagation};
the Lean theorem proves the finite-stage propagation form under
finite-truncation product and defect bounds.  The full statement is
proved here.  See \Cref{subsec:paper1-per-label-coverage}.}
\end{proof}

The operational consequence is that a large enough budget propagates the
membrane.  If \(\widehat K_0\preceqq\Theta_0\) and
\(D_\infty\preceqq\Delta\), then
\(\widehat K_n\preceqq P_\infty(\Theta_0+\Delta)\) for every \(n\).  Thus
the native predictive membrane survives to every accepted stage with a
uniform budget, and this uniform native budget is the input consumed by the
layer-dissolving endpoint of \Cref{sec:dissolution}.

 

\section{Layer-Dissolving Membrane Endpoint}
\label{sec:dissolution}

This section proves the first of the paper's two main results.
\Cref{sec:predictive-membranes} supplied a uniform native budget for
accepted predictive stages, while \Cref{sec:xi-interface} supplied the
adequacy residual as a separately budgeted object.  The endpoint theorem
combines those two inputs with transport into a common dissolving
response space.  The lower-layer obstruction is not declared harmless;
it is priced by native and residual currencies, and that price is
transported into the higher-layer accounting.

At an accepted stage \(j\), the dissolving probe family of
\Cref{sec:xi-interface} is realized as a linear map
\[
  D_j:E_j\to Z_j,
\]
where \(E_j\) is the stage carrier and \(Z_j\) is the stage dissolving
response space.  On the legal quotient at stage \(j\), with \(C_j\) the
stage audit energy restricted to the legal carrier and \(C_j^\dagger\)
its pseudoinverse there, set
\[
  K^L_j=L_jC_j^\dagger L_j^*,\qquad
  K^D_j=D_jC_j^\dagger D_j^*.
\]
Also let
\[
  \Xi_j=\XiRes_{C_j}(D_j\mid L_j),\qquad
  A_j=K^{DL}_j(K^L_j)^\dagger,
\]
where \(K^{DL}_j\) is the dissolving-by-native cross block at stage
\(j\).  The block \(K^L_j\) is the stage-local native self-currency,
\(K^D_j\) is the stage-local dissolving self-currency, \(A_j\) is the
stage-local optimal native explanation map, and \(\Xi_j\) is the
stage-local adequacy residual.  The stagewise Schur identity is
\[
  K^D_j=A_jK^L_jA_j^*+\Xi_j.
\]
This is exactly the adequacy interface theorem of
\Cref{sec:xi-interface} read at stage \(j\).  The predictive currency
\(\widehat K_j\) of \Cref{sec:predictive-membranes} is the routed native
currency in a common native response space; the \(K^L_j\) used here is
the un-routed native block on the stage response space.  Thus the
native-side routing \(R_j\) from \Cref{sec:predictive-membranes} is
distinct from the dissolving-side transport \(S_j\) below, which carries
the dissolving currency into a common target response space.
The superscripts in \(\Theta^Y_j\), \(\Xi^Z_j\), and \(\Theta^D\)
mark, respectively, stage-native, stage-residual, and transported
dissolving target budgets.

\begin{theorem}[Layer-dissolving membrane master theorem]\label{thm:main:layer-dissolving-master}
Assume a formed closure/layer with exact package, null-mode legality,
declared native and dissolving probe families, accepted predictive
transport, all-six records, and nonclaim boundary.  If
\[
  K^L_j\preceqq\Theta^Y_j,
\]
\[
  \Xi_j\preceqq\Xi^Z_j,
\]
and after transport \(S_j:Z_j\to Z\),
\[
  S_j(A_j\Theta^Y_jA_j^*+\Xi^Z_j)S_j^*\preceqq\Theta^D,
\]
then
\[
  S_jK^D_jS_j^*\preceqq\Theta^D.
\]
Consequently, for every layer-dissolving recombination \(z\in Z\),
\[
  \sup_j \Capacity_{C_j}(z^*S_jD_j;E_j)\le z^*\Theta^Dz.
\]
There are no layer-dissolving predictive needles within the declared scope.
\end{theorem}

\begin{proof}
The six named hypotheses specify the scope in which the algebra is
allowed to be read as dissolution.  The exact package and null-mode
legality hypotheses put the stage currencies on the legal quotient; the
declared native and dissolving probe families identify the two probe
blocks; accepted predictive transport supplies the native-side budget;
the all-six records are the transfer records later organized in
\Cref{sec:taxonomy}; and the nonclaim boundary records that the theorem
asserts only this audited endpoint, not a universal singularity-removal
principle.

Start from the stagewise Schur identity
\[
  K^D_j=A_jK^L_jA_j^*+\Xi_j.
\]
Conjugating by the dissolving-side transport \(S_j\) gives
\[
  S_jK^D_jS_j^*
  =S_jA_jK^L_jA_j^*S_j^*+S_j\Xi_jS_j^*.
\]
The native budget \(K^L_j\preceqq\Theta^Y_j\), the residual budget
\(\Xi_j\preceqq\Xi^Z_j\), and Loewner monotonicity under conjugation
therefore imply
\[
  S_jK^D_jS_j^*
  \preceqq
  S_j(A_j\Theta^Y_jA_j^*+\Xi^Z_j)S_j^*.
\]
The transport hypothesis bounds the right-hand side by \(\Theta^D\),
which proves the transported dissolving budget.  The capacity bound and
the no-layer-dissolving-needle conclusion are established below using
this transported bound and the framework currency identity.

For the capacity estimate, fix \(z\in Z\).  The recombination
\(z^*S_jD_j\) is a scalar dissolving probe on the stage carrier
\(E_j\), so the probe-family currency identity from
\Cref{sec:framework} gives
\[
  \Capacity_{C_j}(z^*S_jD_j;E_j)=z^*S_jK^D_jS_j^*z.
\]
Applying the transported dissolving budget yields
\[
  \Capacity_{C_j}(z^*S_jD_j;E_j)\le z^*\Theta^Dz
\]
at every accepted stage.  Taking the supremum over \(j\) gives the
displayed capacity bound.

Finally, a layer-dissolving predictive needle within this paper's scope
would be a dissolving recombination \(z\in Z\) and an accepted stage
\(j\) with strict budget excess
\[
  z^*(S_jK^D_jS_j^*-\Theta^D)z>0.
\]
The transported dissolving budget is exactly the assertion that no such
critical pair exists.  Thus the no-needle conclusion is not an
additional classification theorem; it is the critical-pair translation
of the same Loewner inequality.\footnote{Mechanized in Lean as
\texttt{SixBirdsNeedles.Main.LayerDissolving.layerDissolvingMaster}; the
Lean theorem proves the transport-bound branch under the stagewise
budgets and the transport hypothesis.  The full statement is proved
here.  See \Cref{subsec:paper1-per-label-coverage}.}
\end{proof}

This theorem is the paper's first main dissolution result.  A
lower-layer needle that survives native probes and accepted predictive
certificates is priced rather than dismissed.  The price has two
parts: the native budget \(\Theta^Y_j\) supplied by
\Cref{sec:predictive-membranes}, and the adequacy residual budget
\(\Xi^Z_j\) supplied by \Cref{sec:xi-interface} and developed
quantitatively in \citet{Tsiokos2026XiCompanion}.  After transport by
\(S_j\), the resulting dissolving budget \(\Theta^D\) is uniform across
accepted stages and against arbitrary dissolving recombinations.  The
second main result, the duality-confinement theorem, is the subject of
the next section.

 

\section{Duality-Confinement Membrane Theorem}
\label{sec:duality-confinement}

The previous section priced a layer-dissolving probe family against a
native one through the adequacy interface.  The complementary geometric
setting occurs when the layer carries an involutive duality, such as an
anti-linear or anti-isometric symmetry, and the question is whether
visible object mass is confined to the duality's fixed locus.  The
framework theorem of this section gives conditions under which
anti-invariant membrane collapse implies object-side confinement to
that fixed locus.  Both main results inhabit the same formed-layer theory; they differ
in what is priced and what is confined.

\subsection{Duality-confinement records}

The duality-confinement theorem uses two ledgers.  The first is an
object-side ledger that records how visible objects move under an
involution.  The second is a carrier-side membrane that prices
anti-invariant probe response.  This subsection fixes those records
before the next subsection turns ledger collapse into confinement.
Throughout the section, the superscript ``\(-\)'' marks anti-invariant
carrier-side currencies and budgets.

\begin{definition}[Involutive object ledger]\label{def:main:involutive-ledger}
An involutive object ledger is a tuple
\[
  \mathfrak X=(X,J,\mu,\psi,Y,\mathcal J),
\]
where:
\begin{enumerate}[label=(\roman*)]
\item \(X\) is the visible object space;
\item \(J:X\to X\) is an involution;
\item \(\mu\) is a positive object ledger measure, meaning a positive
  Radon-type measure on \(X\), or a finite positive weight system on
  \(X\);
\item \(Y\) is a real Hilbert response space;
\item \(\mathcal J:Y\to Y\) is a linear isometric involution;
\item \(\psi:X\to Y\) is an equivariant readout:
\[
  \psi(Jx)=\mathcal J\psi(x).
\]
\end{enumerate}
Define
\[
  P_-:=\frac{I-\mathcal J}{2},\qquad \psi_-(x):=P_-\psi(x).
\]
Here \(P_-\) is the spectral projection onto the \(-1\)-eigenspace
of the isometric involution \(\mathcal J\), which is the standard
anti-invariant decomposition for an involution on a Hilbert space
\citep{Conway1990FunctionalAnalysis}.
The anti-invariant object ledger is
\[
  \mathsf A_X:=\int_X \psi_-(x)\psi_-(x)^*\,d\mu(x),
\]
interpreted in the finite-rank or trace-class sense whenever required.
\end{definition}

\begin{definition}[Separating anti-invariant readout]\label{def:main:separating-readout}
The anti-invariant readout separates the fixed locus when
\[
  \psi_-(x)=0 \quad\Longleftrightarrow\quad x\in\Fix(J)
\]
on the visible support of \(\mu\).  It is quantitatively separating if,
for every \(\varepsilon>0\), there exists \(m(\varepsilon)>0\) such that
\[
  \dist(x,\Fix(J))\ge \varepsilon
  \quad\Longrightarrow\quad
  \|\psi_-(x)\|\ge m(\varepsilon).
\]
\end{definition}

\begin{definition}[Carrier-side anti-invariant membrane]\label{def:main:anti-invariant-membrane}
Let \(E\) be an exact packaged carrier with audit \(C\succeq 0\) on its
legal quotient and anti-invariant native probe family
\[
  L^-:E\to Y^-.
\]
The anti-invariant carrier currency is
\[
  \K^-:=L^-C^\dagger (L^-)^*,
\]
with null-mode legality required.  A budget record is a positive
operator \(\Theta^-\) with
\[
  \K^-\preceqq \Theta^-.
\]
A collapse ladder is a family of records
\[
  \K^-_n\preceqq B_n,\qquad \tr B_n\to 0,
\]
or a defected version with \(B_n\) replaced by a bridge budget as
developed in the domination-bridges subsection.
\end{definition}

Together these records separate the two sides of the
duality-confinement argument.  The ledger \(\mathfrak X\) carries the
visible structure and its anti-invariant readout; the separating-readout
condition quantifies how reliably that readout detects departure from
\(\Fix(J)\); and the carrier-side anti-invariant membrane supplies the
framework currency that will be forced to collapse.  The next subsection
turns these records into a confinement conclusion, while the
domination-bridges subsection supplies the comparison from object-side
ledger to carrier-side membrane.

\subsection{Object-side confinement from anti-invariant ledger collapse}

This subsection turns the records of the previous subsection into the
first confinement statements.  The trace identity expresses
\(\mathsf A_X\) as the \(L^2\)-mass of the anti-invariant readout, so
ledger collapse and pointwise vanishing of \(\psi_-\) coincide in the
positive-measure sense.  Combined with the separating-readout condition,
this gives the qualitative and quantitative confinement statements that
anchor the duality-confinement framework.

\begin{lemma}[Trace identity]\label{lem:main:trace-identity}
If \(\mathsf A_X\) is trace-class \citep{Simon2005TraceIdeals}, then
\[
  \tr\mathsf A_X=\int_X \|\psi_-(x)\|^2\,d\mu(x).
\]
Consequently, \(\mathsf A_X=0\) if and only if \(\psi_-(x)=0\) for
\(\mu\)-almost every visible object \(x\).
\end{lemma}

\begin{proof}
The identity follows from the definition of trace for rank-one positive
operators integrated against a positive measure: the trace of
\(\psi_-(x)\psi_-(x)^*\) is \(\|\psi_-(x)\|^2\).  Since the integrand is
nonnegative, its integral vanishes exactly when the integrand vanishes
almost everywhere; and a positive trace-class operator has trace zero
if and only if it is the zero operator
\citep{ReedSimon1980FunctionalAnalysis}.
\end{proof}

\begin{theorem}[Separation implies fixed-locus confinement]\label{thm:main:separation-confinement}
Assume the anti-invariant readout separates the fixed locus.  If
\[
  \mathsf A_X=0,
\]
then
\[
  \mu(X\setminus\Fix(J))=0.
\]
If the ledger is finite with positive weights, every visible object with
positive weight lies in \(\Fix(J)\).
\end{theorem}

\begin{proof}
By \Cref{lem:main:trace-identity}, \(\psi_-(x)=0\) for
\(\mu\)-almost every visible object \(x\).  The separating-readout
condition identifies this zero set with \(\Fix(J)\) on the visible
support of \(\mu\), so the complement of the fixed locus has
\(\mu\)-measure zero.  In the finite positive-weight case, a visible
object outside \(\Fix(J)\) would carry positive mass on the complement,
contradicting that measure-zero conclusion.\footnote{Mechanized in Lean as
\texttt{SixBirdsNeedles.Main.DualityConfinement.separationConfinement};
the Lean theorem proves the finite positive-weight confinement clause
under finite-support separation.  The full statement is proved here.
See \Cref{subsec:paper1-per-label-coverage}.}
\end{proof}

\begin{theorem}[Quantitative confinement]\label{thm:main:quantitative-confinement}
Assume quantitative separation.  Then
\[
  \mu\{x:\dist(x,\Fix(J))\ge\varepsilon\}
  \le
  \frac{\tr\mathsf A_X}{m(\varepsilon)^2}.
\]
\end{theorem}

\begin{proof}
On the set
\(\{x:\dist(x,\Fix(J))\ge\varepsilon\}\), quantitative separation gives
\[
  \|\psi_-(x)\|^2\ge m(\varepsilon)^2.
\]
Applying Markov's inequality \citep{Billingsley1995ProbabilityMeasure}
to the trace identity therefore bounds the
measure of this set by
\(\tr\mathsf A_X/m(\varepsilon)^2\).
\end{proof}

The subsection has converted collapse of the object-side ledger
\(\mathsf A_X\) into confinement to \(\Fix(J)\).  The carrier-side
membrane from the previous subsection has not yet entered the argument;
the next step is to compare it with \(\mathsf A_X\).

\subsection{Domination bridges}

\Cref{thm:main:separation-confinement} obtained confinement from
collapse of the object-side ledger \(\mathsf A_X\), while the
formed-layer framework prices anti-invariant motion on the carrier side
through the membrane currency \(\K^-\).  To use carrier-side membrane
data to drive object-side confinement, the two sides must be related by
a comparison record.  A completed domination bridge is that record;
Douglas domination \citep{Douglas1966Majorization,FillmoreWilliams1971OperatorRanges}
supplies its operator-theoretic foundation; and the
duality-confinement membrane theorem combines bridge collapse with the
trace identity to produce confinement.

\begin{definition}[Completed domination bridge]\label{def:main:domination-bridge}
A completed domination bridge from the object ledger to the carrier
membrane is a record
\[
  \mathsf A_X\preceqq \K^-+E
\]
with \(E\succeq 0\) a declared bridge defect.  Equivalently, in the
exact case \(E=0\), if
\[
  \mathsf A_X=V^*V,\qquad \K^-=W^*W,
\]
then the bridge is exact precisely when there is a contraction \(T\)
such that
\[
  V=TW.
\]
\end{definition}

\begin{theorem}[Douglas domination]\label{thm:main:douglas-domination}
Let \(\mathsf A=V^*V\) and \(\mathsf K=W^*W\) be bounded positive
operators on the same response space.  Then
\[
  \mathsf A\preceqq\mathsf K
\]
if and only if there exists a contraction \(T\) with
\[
  V=TW.
\]
\end{theorem}

\begin{proof}
The reverse direction is immediate.  If \(V=TW\) and \(T\) is a
contraction, then
\[
  V^*V=W^*T^*TW\preceqq W^*W,
\]
because \(T^*T\preceqq I\).

For the forward direction, apply Douglas's factorization lemma to
\(V^*\) and \(W^*\).  The domination
\(V^*V\preceqq W^*W\) is equivalently the range inclusion and uniform
norm bound required by that lemma, so there is a contraction \(T\) with
\(V^*=W^*T^*\), equivalently \(V=TW\).\footnote{Mechanized in Lean as
\texttt{SixBirdsNeedles.Main.DualityConfinement.douglasDominationReverse};
the Lean theorem proves the reverse contraction-to-domination
implication under the stated contraction hypothesis.  The full
statement is proved here.  See \Cref{subsec:paper1-per-label-coverage}.}
\end{proof}

\begin{theorem}[Duality-confinement membrane theorem]\label{thm:main:duality-confinement-master}
Let \(\mathfrak X=(X,J,\mu,\psi,Y,\mathcal J)\) be an involutive object
ledger with separating anti-invariant readout, as in
\Cref{def:main:involutive-ledger,def:main:separating-readout}.  Suppose
there is a sequence of completed domination records
\[
  \mathsf A_X\preceqq B_n \qquad (n=1,2,\ldots)
\]
where \(B_n\succeq 0\) and
\[
  \tr B_n\to 0.
\]
Then
\[
  \mu(X\setminus\Fix(J))=0.
\]
\end{theorem}

\begin{proof}
For every \(n\), the completed domination record gives
\[
  0\preceqq \mathsf A_X\preceqq B_n.
\]
By positivity and trace monotonicity,
\[
  0\le \tr\mathsf A_X\le \tr B_n.
\]
Taking \(n\to\infty\) and using \(\tr B_n\to 0\) gives
\(\tr\mathsf A_X=0\).  Since \(\mathsf A_X\) is positive trace-class,
this implies \(\mathsf A_X=0\).

The separating anti-invariant readout now turns ledger collapse into
object-side confinement.  Applying
\Cref{thm:main:separation-confinement} gives\footnote{Mechanized in Lean as
\texttt{SixBirdsNeedles.Main.DualityConfinement.dualityConfinementMaster};
the Lean theorem proves a finite-rank single-bridge version with
positive-weight confinement.  The full statement is proved here.  See
\Cref{subsec:paper1-per-label-coverage}.}
\[
  \mu(X\setminus\Fix(J))=0.
\]
\end{proof}

This subsection delivers the master theorem of
\Cref{sec:duality-confinement}.  The two-ledger records of the first
subsection, the trace-identity reasoning of the second, and the
comparison apparatus above together give the conditional confinement
promised in the section opener.  The remaining subsections supply the
finite-ledger and defect-budget machinery needed to use that theorem.

\subsection{Exhaustive finite ledgers}

\Cref{thm:main:duality-confinement-master} requires a sequence of
completed domination records that drives the object ledger to zero.  In
many applications the completed ledger is not accessed directly; one
has finite-window or moving approximations instead.  Exhaustive finite
ledgers package those approximations so that vanishing exhaustivity
supplies the input required by the master theorem.

\begin{definition}[Exhaustive moving ledger]\label{def:main:exhaustive-moving-ledger}
A finite-window or moving object ledger \(\mathsf A_{X,n}\) is
exhaustive for a completed ledger \(\mathsf A_X\) if there are
inclusions or transports \(\iota_n\) and tail operators
\(T_n\succeq 0\) such that
\[
  \mathsf A_X\preceqq \iota_n\mathsf A_{X,n}\iota_n^*+T_n.
\]
It is vanishing-exhaustive under bounds
\(\mathsf A_{X,n}\preceqq B_n\) if
\[
  \tr(\iota_nB_n\iota_n^*)+\tr T_n\to 0.
\]
\end{definition}

\begin{theorem}[Exhaustive ledger squeeze]\label{thm:main:exhaustive-ledger-squeeze}
If \(\mathsf A_{X,n}\preceqq B_n\) and the moving ledger is
vanishing-exhaustive for \(\mathsf A_X\), then
\[
  \mathsf A_X=0.
\]
Consequently, under readout separation, visible object mass is confined
to \(\Fix(J)\).
\end{theorem}

\begin{proof}
By exhaustivity and the bounds \(\mathsf A_{X,n}\preceqq B_n\),
\[
  0\preceqq\mathsf A_X
  \preceqq \iota_nB_n\iota_n^*+T_n.
\]
Taking traces gives
\[
  0\le \tr\mathsf A_X
  \le \tr(\iota_nB_n\iota_n^*)+\tr T_n.
\]
The right-hand side tends to \(0\) by vanishing exhaustivity, so
\(\tr\mathsf A_X=0\).  Since \(\mathsf A_X\) is positive trace-class,
this implies \(\mathsf A_X=0\).  Under readout separation, the
confinement consequence follows from
\Cref{thm:main:separation-confinement}.
\end{proof}

The exhaustive-ledger squeeze packages the typical application pattern:
finite or moving approximations dominate the full ledger with shrinking
tails.  The remaining subsection handles the case where the domination
records themselves come with declared defects.

\subsection{Defected obstruction budgets}

The master theorem and the exhaustive squeeze use completed domination
records with no residual defect.  In practice, both the bridge from
object ledger to carrier membrane and the source-side membrane estimate
may carry declared positive defects.  This subsection packages the
defected obstruction budget and records the scalar optimization that
turns those defects into a single sufficient sublinear condition.

\begin{definition}[General obstruction budget]\label{def:main:obstruction-budget}
A defected duality-confinement record has the form
\[
  \mathsf A_X
  \preceqq
  (1+t)(\K^-_n+E_{\mathrm{src},n})
  +(1+t^{-1})E_{\mathrm{br},n}
\]
for \(t>0\), where \(E_{\mathrm{br},n}\) is bridge/domination defect
and \(E_{\mathrm{src},n}\) is source/membrane defect.  In a collapse
ladder one often has
\[
  \K^-_n\preceqq \Lambda_n^{-1}\Theta_0^-.
\]
Then the budget is
\[
  B_n(t)
  =(1+t)(\Lambda_n^{-1}\Theta_0^-+E_{\mathrm{src},n})
   +(1+t^{-1})E_{\mathrm{br},n}.
\]
\end{definition}

This is the operator Cauchy-Schwarz or Young splitting of the cross
term \citep{BhatiaDavis1995CauchySchwarz,Ando1995MatrixYoung} that
appears when both the source membrane and the domination bridge carry
positive defects; the parameter \(t\) balances the two defect
contributions.

\begin{proposition}[Optimized scalar trace budget]\label{prop:main:optimized-trace-budget}
Let
\[
  a_n=\tr(\Lambda_n^{-1}\Theta_0^-+E_{\mathrm{src},n}),
  \qquad
  b_n=\tr E_{\mathrm{br},n}.
\]
Then
\[
  \inf_{t>0}\tr B_n(t)=(\sqrt{a_n}+\sqrt{b_n})^2.
\]
Thus a sufficient scalar condition for exact confinement by the
exhaustive-ledger squeeze is
\[
  a_n\to 0,\qquad b_n\to 0.
\]
\end{proposition}

\begin{proof}
Taking traces in the definition of \(B_n(t)\) gives
\[
  \tr B_n(t)=(1+t)a_n+(1+t^{-1})b_n.
\]
For \(a_n,b_n>0\), differentiating in \(t>0\) gives the minimizer
\[
  t=\sqrt{b_n/a_n},
\]
and substitution gives
\[
  \inf_{t>0}\tr B_n(t)=(\sqrt{a_n}+\sqrt{b_n})^2.
\]
If one of \(a_n,b_n\) is zero, the same value follows by continuity, or
by taking the corresponding endpoint limit in \(t\).

If \(a_n\to0\) and \(b_n\to0\), then the optimized traces tend to zero.
Choosing \(t_n\) at, or arbitrarily near, the scalar infimum produces
budgets \(B_n(t_n)\) with \(\tr B_n(t_n)\to0\).  The
exhaustive-ledger squeeze of
\Cref{thm:main:exhaustive-ledger-squeeze} then gives exact
confinement.
\end{proof}

This subsection closes the framework apparatus for
duality-confinement.  The five subsections have isolated the involutive
object ledger and carrier-side anti-invariant membrane records, the
trace identity connecting anti-invariant ledger collapse to fixed-locus
confinement, the domination bridges and master theorem, the
exhaustive-finite-ledger squeeze, and finally the defected obstruction
budget that admits declared bridge and source defects.  Falsifiability
is therefore record-level: an instance that fails to achieve confinement
must fail by lacking, weakening, or defecting one of these records, in a
way exposed by the layer-dissolving and duality-confinement endpoints.
This concludes the body's treatment of the duality-confinement
mechanism; the next section turns to the critical-pair-completion and
strict-extension records that the membrane theorems rely on, and
\Cref{sec:taxonomy} organizes the all-six taxonomy that scopes when a
record is admissible.

 

\section{Critical-Pair Completion and Strict Extension}
\label{sec:critical-pairs}

The layer-dissolving endpoint, \Cref{thm:main:layer-dissolving-master},
and the duality-confinement master theorem,
\Cref{thm:main:duality-confinement-master}, both use critical-pair
language---adapted here from the rewriting-systems usage of
\citet{KnuthBendix1970SimpleWord} and the modern exposition of
\citet{BaaderNipkow1998TermRewriting}---to translate failure of a
Loewner budget into a witness pair.
This section makes that translation explicit on the native and
residual sides of the adequacy interface, records what counts as a
lawful repair, and proves the strict-extension contraction that
explains why adding genuinely new native probes can only decrease the
adequacy residual.

A native-side budget failure has the following form.  If \(K\) denotes
one of the native currencies used earlier in the paper---for example
the packaged native currency \(\K_{C,\Gamma}(L)\) of \Cref{sec:framework}
or the predictive currency \(\widehat K_j\) of
\Cref{sec:predictive-membranes}---then failure of
\[
  K\preceqq\Theta
\]
is equivalent to the existence of a response vector \(y\) with
\[
  y^*(K-\Theta)y>0.
\]
The vector \(y\), together with the budget it violates, is the
native-side critical pair.  Similarly, if \(\Xi\) denotes the adequacy
residual of \Cref{sec:xi-interface} and \(\Xi_Z\) a declared residual
budget, failure of
\[
  \Xi\preceqq\Xi_Z
\]
is equivalent to the existence of a dissolving response vector \(z\)
with
\[
  z^*(\Xi-\Xi_Z)z>0.
\]
This pair \((\Xi,z)\) is the residual-side critical pair.  Thus a
critical pair is not extra structure; it is the vector witness supplied
by Loewner separation when a declared budget fails.

The framework allows only certain repairs to such a failure.  A repair
is lawful if it excludes the offending recombination by a declared
nonclaim, weakens the claim by raising the relevant budget, strengthens
or rewrites the carrier and audit so that the legal quotient has
changed, adds native probes that see the residual direction, or
completes the missing protocol or transport family and then retests the
same block.  These options are deliberately narrow: they change either
the scope of the claim, the budget being asserted, the legal carrier on
which the currency is computed, the probe family available to explain
the residual, or the transfer record that transports the budget.  By
contrast, same-package repetition saturates.  Repeating the same
Schur-block repair with no new probe content, no new audit, and no new
transport data does not change the residual and therefore cannot remove
the critical pair.

\begin{theorem}[Strict extension contraction]\label{thm:main:strict-extension-contraction}
If a strict native extension adds probes
\[
  L^+=\begin{bmatrix}L\\M\end{bmatrix},
\]
then
\[
  \XiRes_C(D\mid L^+)
  =
  \XiRes_C(D\mid L)
  -K_{DM\mid L}K_{MM\mid L}^\dagger K_{MD\mid L}
  \preceqq \XiRes_C(D\mid L).
\]
Strict contraction occurs exactly in the directions seen by the new probes.
\end{theorem}

\begin{proof}
The theorem is the Schur-update form of the adequacy residual,
in the lineage of the Crabtree--Haynsworth quotient identity
\citep{CrabtreeHaynsworth1969Quotient} and the operator-theoretic
treatment of Schur complements in \citep{Ando2005SchurOperator}.  Passing
from \(L\) to
\[
  L^+=\begin{bmatrix}L\\M\end{bmatrix}
\]
enlarges the native block in the joint dissolving-native currency.  The
new block \(M\) has a conditional residual \(K_{MM\mid L}\), obtained by
subtracting the part already explained by \(L\), and cross terms
\(K_{DM\mid L}\) and \(K_{MD\mid L}\) against the dissolving block.  The
chain rule for adequacy residuals gives
\[
  \XiRes_C(D\mid L^+)
  =
  \XiRes_C(D\mid L)
  -K_{DM\mid L}K_{MM\mid L}^\dagger K_{MD\mid L},
\]
and the full residual chain-rule calculus is developed in
\citet{Tsiokos2026XiCompanion}.  In this paper the displayed identity is
used only as the strict-extension update governing lawful repairs by
new native probes.

The Loewner contraction follows because the Schur correction
\[
  K_{DM\mid L}K_{MM\mid L}^\dagger K_{MD\mid L}
\]
is positive semidefinite.  Subtracting a positive semidefinite
correction from \(\XiRes_C(D\mid L)\) can only decrease the residual in
the Loewner order, giving
\[
  \XiRes_C(D\mid L^+)\preceqq\XiRes_C(D\mid L).
\]
The strict-contraction-along-seen-directions content follows from
positive semidefiniteness of the Schur correction.

Finally, the correction is nonzero exactly where the newly added probes
have conditional content beyond what the old native family already
explained.  If \(M\) sees no additional conditional direction, then the
Schur correction vanishes and the residual is unchanged.  If it sees a
new direction coupled to the dissolving block, that direction is
subtracted from the residual budget.  This is the precise sense in
which strict extension contracts the adequacy residual only along
directions seen by the new probes.\footnote{Mechanized in Lean as
\texttt{SixBirdsNeedles.Main.CriticalPair.strictExtensionContraction};
the Lean theorem proves the Loewner inequality under elevated
Schur-update hypotheses.  The full statement is proved here.  See
\Cref{subsec:paper1-per-label-coverage}.}
\end{proof}

Critical pairs are therefore the operational form that Loewner failure
takes inside the formed-layer framework.  Lawful repair is constrained
to the five mechanisms above, and strict extension contraction
quantifies the repair obtained by adding native probes: the adequacy
residual monotonically shrinks as genuinely new native information is
made available.  \Cref{sec:taxonomy} records the all-six discipline
that decides when such a repair record is admissible as a lawful
transfer rather than an informal analogy.

 

\section{All-Six Records and Adequacy-Aware Transfer}
\label{sec:taxonomy}

The two main results---the layer-dissolving endpoint of
\Cref{sec:dissolution} and the duality-confinement master theorem of
\Cref{sec:duality-confinement}---together with the strict-extension
contraction of \Cref{sec:critical-pairs}, all depend on the records
that justify their hypotheses.  This section organizes those records
into a six-role checklist, the \emph{all-six taxonomy}.  The taxonomy
decides whether a proposed transfer is a lawful membrane transfer
rather than an analogy between layers, and the adequacy-aware transfer
condition records the residual budgets that must move with it.

The six primitive roles are
\[
  P_1:\text{ rewrite/gauge},\quad
  P_2:\text{ feasibility/channel},\quad
  P_3:\text{ route/holonomy},
\]
\[
  P_4:\text{ staging/refinement},\quad
  P_5:\text{ packaging/canonicalization},\quad
  P_6:\text{ audit/currency}.
\]
Here \(P_1\) is the formal rewrite or change of presentation, such as
replacing a lower-layer obstruction statement by a higher-layer
residual statement.  \(P_2\) records admissibility of the maps and
constraints being used, including null-legality of native and dissolving
probes and well-defined Loewner inequalities.  \(P_3\) records how data
moves through the required transport maps, such as the dissolving-side
transport \(S_j\) in \Cref{sec:dissolution}.  \(P_4\) records the finite
stage, refinement, or accepted index at which the claim is made, as in
the accepted predictive stages of \Cref{sec:predictive-membranes}.
\(P_5\) records the canonical package or quotient in which the data is
represented, including the legal carrier \(E_C\) and the exact package
\(\Gamma\) of \Cref{sec:framework}.  \(P_6\) records the currency,
budget, or checkable accounting that validates the transfer, such as a
Loewner budget on native currency, dissolving currency, or adequacy
residual.

Omitting any of these six channels yields a finite overclaim
countermodel: the formal rewrite or budget claim becomes either
ill-defined or unaudited, and the framework treats such an omission as
a failure of the record rather than a failure of the underlying
mathematics.  This is the role of the taxonomy.  It is not meant to
name every possible obstruction; it identifies which records must be
present before the endpoint theorems can be read as transfer across
layers.

Membranes may appear in several strengths.  An \emph{intrinsic}
membrane is stated directly on the target object, without auxiliary
choices; a \emph{sectioned} membrane depends on a chosen section or
presentation; a \emph{bundle/fiber} membrane is stated fiberwise over a
base; and a \emph{protocol-indexed} membrane is indexed by an accepted
process or refinement stage.  Intrinsic form is the strongest.  A
weaker membrane promotes to intrinsic only through faithful
reconstruction, which supplies the \(P_5\) and \(P_1\) records; full
block currency control, which supplies the \(P_6\) audit; and defect
accounting, again a \(P_6\) requirement and possibly a defected budget
of the kind used in \Cref{sec:duality-confinement}.

\begin{definition}[Lawful transfer record]\label{def:main:lawful-transfer}
A lawful transfer record between two exact packaged carriers
\(E_{\mathrm{src}}\) and \(E_{\mathrm{tgt}}\), with audit operators
\(C_{\mathrm{src}},C_{\mathrm{tgt}}\), native probe families
\(L_{\mathrm{src}},L_{\mathrm{tgt}}\), dissolving probe families
\(D_{\mathrm{src}},D_{\mathrm{tgt}}\), and adequacy residuals
\(\XiRes^{\mathrm{src}},\XiRes^{\mathrm{tgt}}\), consists of four
transport maps
\[
  T_E:E_{\mathrm{src}}\to E_{\mathrm{tgt}},\quad
  T_L:Y_{\mathrm{src}}\to Y_{\mathrm{tgt}},\quad
  T_D:Z_{\mathrm{src}}\to Z_{\mathrm{tgt}},\quad
  T_\Xi:Z_{\mathrm{src}}\to Z_{\mathrm{tgt}},
\]
each declaring its role under \(P_1,\ldots,P_6\), together with the
currency-block compatibility data declared in the next two definitions.
The transport \(T_\Xi\) carries the adequacy-residual budget on the
dissolving response space.
\end{definition}

\begin{definition}[Exact bridge]\label{def:main:exact-bridge}
A lawful transfer record is an exact bridge when the transport maps
deliver currency-block congruences on the source and target legal
quotients:
\[
  T_L^*K^{\mathrm{tgt}}_{LL}T_L=K^{\mathrm{src}}_{LL},
  \qquad
  T_D^*K^{\mathrm{tgt}}_{DD}T_D=K^{\mathrm{src}}_{DD},
\]
and the residual-transport map \(T_\Xi\) satisfies the adequacy-residual
congruence
\[
  T_\Xi^*\,\XiRes_{C_{\mathrm{tgt}}}(D_{\mathrm{tgt}}\mid L_{\mathrm{tgt}})\,
  T_\Xi
  =
  \XiRes_{C_{\mathrm{src}}}(D_{\mathrm{src}}\mid L_{\mathrm{src}}).
\]
\end{definition}

\begin{definition}[Defective bridge]\label{def:main:defective-bridge}
A lawful transfer record is a defective bridge with defect operator
\(\Delta\succeq0\) on the target dissolving block when the transport
maps deliver
\[
  T_L^*K^{\mathrm{tgt}}_{LL}T_L=K^{\mathrm{src}}_{LL},
\]
together with the currency-block Loewner domination
\[
  K^{\mathrm{tgt}}_{DD}\preceqq
  T_D\,K^{\mathrm{src}}_{DD}\,T_D^*+\Delta
\]
on the target dissolving block.
\end{definition}

\begin{proposition}[Adequacy-aware transfer]\label{thm:main:adequacy-aware-transfer}
A lawful transfer of a layer-dissolving membrane must transport energy,
native probes, dissolving probes, and adequacy residuals.  Exact bridges
transfer the native and dissolving budgets by congruence.  Defective
bridges produce explicit residual currencies; these must be paid in the
target budget.
\end{proposition}

\begin{proof}
A lawful transfer record, in the sense of
\Cref{def:main:lawful-transfer}, declares each of \(P_1,\ldots,P_6\)
for the transported energy, native probes, dissolving probes, and
adequacy residuals.  If one of these roles is omitted, then a channel
in the all-six checklist is undeclared, and the proposed transfer is
not eligible to carry a membrane.  This proves the first clause of the
statement.

For the congruence clause, suppose the transfer record is an exact
bridge as in \Cref{def:main:exact-bridge}.  Then the currency-block
congruences
\[
  T_L^*K^{\mathrm{tgt}}_{LL}T_L=K^{\mathrm{src}}_{LL},
  \qquad
  T_D^*K^{\mathrm{tgt}}_{DD}T_D=K^{\mathrm{src}}_{DD},
\]
and the adequacy-residual congruence
\[
  T_\Xi^*\XiRes_{C_{\mathrm{tgt}}}(D_{\mathrm{tgt}}\mid L_{\mathrm{tgt}})
  T_\Xi
  =
  \XiRes_{C_{\mathrm{src}}}(D_{\mathrm{src}}\mid L_{\mathrm{src}})
\]
are exactly the budget-level assertion that native, dissolving, and
adequacy-residual budgets transfer by congruence.  No further
calculation is required: the exact-bridge definition is the congruence
statement.

For the residual-payment clause, suppose the transfer record is a
defective bridge with defect operator \(\Delta\succeq0\) on the target
dissolving block, as in \Cref{def:main:defective-bridge}.  Then the
domination
\[
  K^{\mathrm{tgt}}_{DD}
    \preceqq T_DK^{\mathrm{src}}_{DD}T_D^*+\Delta,
\]
is the Loewner statement that the target dissolving currency is bounded
by the transported source dissolving currency plus the explicit
positive defect \(\Delta\).  The defect operator \(\Delta\) is therefore
the residual currency that must be paid in the target adequacy budget;
the budget-aware transfer used in \Cref{sec:dissolution} routes it
through the dissolving budget on the target side.\footnote{Tracked in the Lean formalization as
\texttt{SixBirdsNeedles.Main.AllSix.adequacyAwareTransfer}: a conditional
schema over the lawful-transfer record fields rather than an independent
derivation of the quantitative all-six transfer conditions, which are
proved here.  See \Cref{subsec:paper1-per-label-coverage}.}
\end{proof}

The quantitative all-six conditions are made explicit by the per-budget
statements of \Cref{sec:dissolution} and \Cref{sec:duality-confinement},
which respectively supply the optimized scalar trace budget and the
layer-dissolving budget consumed in the congruence and residual-payment
clauses above.

The all-six discipline and adequacy-aware transfer together fix what
counts here as lawful emergence: a dissolution claim is admissible only
when its records cover all six roles and the transport carries adequacy
data with declared budgets.  This taxonomy does not classify all
needles or all possible repairs.  It records the roles required for the
instances treated in this paper, so that the conclusion is a checked
transfer from lower-layer obstruction to higher-layer dissolution.  The
remaining sections discuss the scope of that transfer and the limits of
the present framework.

 

\section{Discussion}
\label{sec:discussion}

The paper has proved two framework-bound dissolution mechanisms.  In
the formed-layer membrane setting, where the legal quotient, packaged
native currency, adequacy residual, predictive native membrane, and
dissolving transport records are all present, the layer-dissolving
endpoint converts a bounded native budget and a bounded adequacy
residual into a uniform dissolving budget.  Within that declared scope,
there are no layer-dissolving predictive needles.  The complementary
duality-confinement theorem turns anti-invariant ledger collapse into
fixed-locus confinement under separating readouts.  These are
formed-layer dissolution mechanisms, not universal
singularity-elimination results.

The local mathematics used by the framework is classical.  Positive
semidefinite quadratic forms, Loewner order, Moore--Penrose
pseudoinverses, Schur complements, and Douglas factorization enter in
standard ways, and the paper does not claim novelty for those
constructions.  The contribution is the synchronization of these
constructions across a passage between layers.  The quotient used by
the probes must be the quotient used by the budget; the native currency
must be the currency spent by the predictive membrane; the adequacy
residual must be the part of the dissolving currency that native probes
do not explain; and the endpoint theorem must spend exactly those
currencies.  This is the framework-as-discipline claim: the same
operator tools become a dissolution theorem only when their roles are
kept aligned.

The two main results are not redundant.  The layer-dissolving
endpoint requires a dissolving response space, stagewise residual
budgets, and transport maps \(S_j\) into a common target; it prices the
dissolving block from above.  The duality-confinement theorem requires
an involutive object ledger, a separating anti-invariant readout, and
domination records that force the anti-invariant ledger to collapse; it
confines visible object mass to the fixed locus of the involution.  The
two endpoints share the legal-quotient discipline, the formed-layer
budget language, the all-six taxonomy, and the adequacy-aware transfer
condition, but their hypothesis classes are not nested.  A problem with
no involutive duality has no duality-confinement instance, while a
problem with no transport into a common dissolving response space has
no layer-dissolving endpoint instance.

The separation between this paper and the residual-calculus development
of \citet{Tsiokos2026XiCompanion} is also part of the architecture.  The
present paper owns the emergence-and-dissolution narrative: it states the
adequacy interface theorem (\Cref{thm:main:adequacy-interface}) as the
gate through which native and residual budgets enter the endpoint.  The
adequacy-residual calculus---projection, optimality, exactness, chain
rule, data-processing, and witness theorems---is developed in
\citet{Tsiokos2026XiCompanion}.  Keeping the developments separate lets
the present paper state the dissolution mechanisms without forcing every
reader through the full Schur-complement calculus, while leaving the
residual-calculus development self-contained for readers who need that
theory in its own right.

The scope fence is correspondingly explicit.  The paper does not claim
that emergence dissolves every mathematical singularity or obstruction,
does not give a complete classification of needles or repairs, and does
not depend on any particular application instance.  It also does not
claim novelty for the classical operator-theoretic constructions it
uses.  The result is a framework theorem: needle-class obstructions
satisfying the stated predictive, adequacy, endpoint,
duality-confinement, and all-six records can be dissolved in the
formed-layer membrane setting.  Future application papers may aim this
machinery at specific external problems, but each such use would need
its own framework-specific records, scope fence, and formalization
disclosure.

Several directions remain open.  The strict-extension contraction of
\Cref{thm:main:strict-extension-contraction} suggests a quantitative
theory of probe-family enlargement, with the residual chain rule of
\citet{Tsiokos2026XiCompanion} as the natural tool.  The defected obstruction budget
machinery in \Cref{sec:duality-confinement} may admit sharper tradeoffs
between source and bridge defects under stronger structural assumptions
on the involution.  Finally, the all-six taxonomy and
adequacy-aware-transfer discipline may apply to transfer settings beyond
layer-dissolving emergence, but such uses would require their own
audited instances and are outside the present paper.

 

\section{Conclusion}
\label{sec:conclusion}

The paper's main claim has been proved in framework-bound form.  In
the formed-layer membrane setting, emergence can dissolve needle-class
lower-layer obstructions when the legal quotient, native currency,
adequacy residual, predictive membrane bounds, dissolving transport
records, and all-six admissibility discipline are all present.  Two
complementary main results establish this conclusion: the layer-dissolving
endpoint and the duality-confinement membrane theorem, supported by
the adequacy interface, the strict-extension contraction, and the
all-six taxonomy.

The dissolution claim is therefore instance-specific.  It is not a
universal singularity-elimination theorem, not a complete
classification of needles or repairs, and not an application proof for
any named external problem in this submission.  Future work that aims
the framework at a specific external problem must supply its own
application-specific records, scope fence, and formalization disclosure.

The present paper uses the adequacy residual only as a gate, through
the adequacy interface theorem (\Cref{thm:main:adequacy-interface}).  The full residual
calculus, including projection, optimality, exactness, chain-rule,
data-processing, and witness results, is developed in
\citet{Tsiokos2026XiCompanion}.

Several concrete directions remain.  The all-six discipline may extend
to transfer settings beyond layer-dissolving emergence; the
defect-budget tradeoffs in duality-confinement may sharpen under
stronger involution hypotheses; and explicit audited application
instances may exercise every record required by the framework.  These
are further instances of the rule already used here: emergence
dissolves needles only when the layer change is paid for by the
correct records.

 
\appendix
\section{Lean Formalization Disclosure}
\label{app:formalization}

\subsection{Trust base and methodology}
\label{subsec:paper1-trust-methodology}

The Lean artifacts cited in this appendix were checked with Lean 4.28.0 \citep{deMouraUllrich2021Lean4,MathlibCommunity2020Lean}.  For every cited theorem entry, the audited axiom closure is contained in the standard trust base: \texttt{propext}, \texttt{Classical.choice}, and \texttt{Quot.sound}.  The project discipline forbids \texttt{sorry}, project-local \texttt{axiom}, and \texttt{opaque} declarations in the needles Lean corpus; \texttt{lean/manifests/trust\_base.txt} is the locked allowlist used by the manifest checker when auditing theorem dependency closures.  Coverage is recorded per row as \emph{theorem}, \emph{partial}, \emph{definition}, or \emph{not mechanized}; these labels are recorded in the per-label coverage table below.

\subsection{Mechanization disclosure tiers}
\label{subsec:paper1-mechanization-disclosure-tiers}

Body footnotes use a three-tier disclosure discipline keyed to the alignment column of the per-label coverage table.

\begin{table}[H]
  \centering
  \small
  \begin{tabularx}{\linewidth}{>{\raggedright\arraybackslash}p{0.25\linewidth} >{\raggedright\arraybackslash}p{0.11\linewidth} >{\raggedright\arraybackslash}X}
    \toprule
    Alignment & Tier label & Body footnote wording template \\
    \midrule
    faithful & T1 & Verified in Lean as \(\langle\)module\(\rangle\). \\
    narrowed-surrogate / needs-strengthening & T2 & Mechanized in Lean as \(\langle\)module\(\rangle\); proves \(\langle\)sub-claim\(\rangle\) under \(\langle\)elevation\(\rangle\); full statement paper-side. \\
    projection-packaged & T3 & Tracked in the Lean formalization as \(\langle\)module\(\rangle\): conditional schema over \(\langle\)record fields\(\rangle\), not an independent derivation. \\
    \bottomrule
  \end{tabularx}
  \caption{Body footnote disclosure tiers.}
  \label{tab:paper1-mechanization-disclosure-tiers}
\end{table}

This tier discipline is the canonical wording rule for body footnotes: a reader can cross-check each body footnote against the alignment column of the coverage table below.

\subsection{Standing narrowings and elevated hypotheses}
\label{subsec:paper1-standing-narrowings}

Several Lean fragments use the concrete diagonal pseudoinverse \texttt{diagonalPseudoInverse C} as the finite-matrix representative of the paper's Moore--Penrose inverse \(C^\dagger\).  This is paper-faithful on the legal quotient under the corresponding strict-positivity support hypotheses, but the full quotient-level pseudoinverse theory is not mechanized in this codebase.  When the body invokes \(C^\dagger\), the per-label disclosure below records whether Lean proves the stated clause directly or proves a finite surrogate under stricter positivity.

Many partial-status theorems also carry elevated algebraic hypotheses, such as transposition symmetry, Moore--Penrose-on-support identities, or explicit Schur-update assumptions.  The elevation is documented per row in the coverage table; the paper supplies the ordinary mathematical proof or explanation for the remaining clause.  These elevations should be stated in mathematical language when relevant and should not be hidden behind manifest terminology.

The projection-packaged status flags statements whose Lean form retrieves data from a proof-carrying record \citep{Necula1997ProofCarryingCode} or conditional schema rather than deriving the paper theorem independently.  Such rows are traceability coverage, not substantive Lean proof coverage.  The body relies on the mathematical statement, while the coverage table records the traceability role.

\subsection{Per-label coverage table}
\label{subsec:paper1-per-label-coverage}

Tables~\ref{tab:paper1-lean-coverage} and~\ref{tab:paper1-lean-coverage-continued}
record the per-row Lean coverage, alignment, and short disclosure
note for each statements-of-record row.
Table~\ref{tab:paper1-nonmechanized-coverage} lists the rows retained
with no Lean entry.

\begin{table}[!ht]
  \centering
  \scriptsize
  \setlength{\tabcolsep}{2.5pt}
  \renewcommand{\arraystretch}{0.92}
  \begin{tabularx}{\linewidth}{>{\raggedright\arraybackslash}p{0.38\linewidth} >{\raggedright\arraybackslash}p{0.11\linewidth} >{\raggedright\arraybackslash}p{0.15\linewidth} >{\raggedright\arraybackslash}X}
    \toprule
    Paper label & Lean coverage & Alignment & Disclosure note \\
    \midrule
    \path{def:main:packaged-currency} & definition & faithful & Realized in Lean as the packaged native currency definition. \\
    \path{thm:main:probe-family-currency} & partial & narrowed-surrogate & Lean proves the point-evaluation matrix identity under stricter positivity; the variational capacity characterization is established in the paper. \\
    \path{thm:main:minimum-spend-duality} & partial & needs-strengthening & Lean proves the off-range infinite-cost branch; the finite-range branch is supplied as paper-side proof. \\
    \path{thm:main:adequacy-interface} & partial & needs-strengthening & Lean proves the Schur identity under explicit algebraic hypotheses; Loewner monotonicity and exact-factorization equivalence are paper-side. \\
    \path{thm:main:predictive-native-membrane} & theorem & faithful & Lean proves the faithful finite predictive-membrane statement. \\
    \path{thm:main:summable-predictive-propagation} & partial & narrowed-surrogate & Lean proves the finite-stage propagation form; infinite convergence remains a paper-side limiting statement. \\
    \path{thm:main:layer-dissolving-master} & partial & needs-strengthening & Lean proves the transport-bound branch; capacity-bound and no-needle conclusions are paper-side. \\
    \path{def:main:involutive-ledger} & partial & narrowed-surrogate & Lean records the finite-rank rational-matrix object ledger; measure-theoretic generality is paper-side. \\
    \path{def:main:separating-readout} & partial & narrowed-surrogate & Lean records finite-support separation and a rational lower-bound variant. \\
    \path{def:main:anti-invariant-membrane} & partial & narrowed-surrogate & Lean records the finite-rank carrier currency with diagonal pseudoinverse infrastructure. \\
    \path{lem:main:trace-identity} & partial & needs-strengthening & Lean proves the finite-sum trace equality; the almost-everywhere consequence is paper-side. \\
    \bottomrule
  \end{tabularx}
  \caption{Manifest coverage at commit \texttt{3bbe9f9}, part 1.}
  \label{tab:paper1-lean-coverage}
\end{table}

\begin{table}[!ht]
  \centering
  \scriptsize
  \setlength{\tabcolsep}{2.5pt}
  \renewcommand{\arraystretch}{0.92}
  \begin{tabularx}{\linewidth}{>{\raggedright\arraybackslash}p{0.38\linewidth} >{\raggedright\arraybackslash}p{0.11\linewidth} >{\raggedright\arraybackslash}p{0.15\linewidth} >{\raggedright\arraybackslash}X}
    \toprule
    Paper label & Lean coverage & Alignment & Disclosure note \\
    \midrule
    \path{thm:main:separation-confinement} & partial & needs-strengthening & Lean proves the finite positive-weight confinement clause; the measure-zero formulation is paper-side. \\
    \path{thm:main:quantitative-confinement} & partial & narrowed-surrogate & Lean proves the finite-rank rational Markov-style bound. \\
    \path{def:main:domination-bridge} & partial & narrowed-surrogate & Lean records the finite-rank domination bridge predicate and related factorization predicates. \\
    \path{thm:main:douglas-domination} & partial & needs-strengthening & Lean proves only the reverse contraction-to-domination direction; the forward Douglas factorization is paper-side. \\
    \path{thm:main:duality-confinement-master} & partial & narrowed-surrogate & Lean proves a finite-rank single-bridge version with positive-weight confinement. \\
    \path{def:main:exhaustive-moving-ledger} & partial & narrowed-surrogate & Lean records a single finite exhaustive ledger relation rather than a sequence with limits. \\
    \path{thm:main:exhaustive-ledger-squeeze} & partial & narrowed-surrogate & Lean proves a single-record finite-rank squeeze; sequence and limit language remains paper-side. \\
    \path{def:main:obstruction-budget} & partial & narrowed-surrogate & Lean records the finite-rank per-parameter obstruction budget. \\
    \path{prop:main:optimized-trace-budget} & partial & needs-strengthening & Lean proves the per-\(t\) trace formula; the closed-form infimum and convergence clause are paper-side. \\
    \path{thm:main:strict-extension-contraction} & partial & needs-strengthening & Lean proves the Loewner inequality from elevated Schur-update hypotheses; strict seen-direction content is paper-side. \\
    \path{thm:main:adequacy-aware-transfer} & partial & projection-packaged & Lean records data shape and certificate fields only; quantitative all-six transfer is paper-side. \\
    \bottomrule
  \end{tabularx}
  \caption{Manifest coverage at commit \texttt{3bbe9f9}, part 2.}
  \label{tab:paper1-lean-coverage-continued}
\end{table}

\begin{table}[!ht]
  \centering
  \scriptsize
  \setlength{\tabcolsep}{2.5pt}
  \begin{tabularx}{\linewidth}{>{\raggedright\arraybackslash}p{0.40\linewidth} >{\raggedright\arraybackslash}p{0.14\linewidth} >{\raggedright\arraybackslash}X}
    \toprule
    Paper label & Destination & Disclosure note \\
    \midrule
    \path{rmk:main:blind-spot-interpretation} & appendix & No Lean entry; interpretive material is handled in prose, with the body-facing blind-spot calculus developed in \citet{Tsiokos2026XiCompanion}. \\
    \path{warn:main:defect-paid-version} & appendix & No Lean entry; this variant warning is retained as appendix support and is not part of the body theorem spine. \\
    \bottomrule
  \end{tabularx}
  \caption{Retained rows with no Lean entry.}
  \label{tab:paper1-nonmechanized-coverage}
\end{table}

\subsection{Lean declaration map}
\label{subsec:paper1-lean-declaration-map}

Table~\ref{tab:paper1-lean-declaration-map} maps each paper label to
its full Lean declaration path; the per-row coverage notes above use
short prose disclosures and refer to this table for the canonical
identifiers.

\begin{table}[!ht]
  \centering
  \scriptsize
  \setlength{\tabcolsep}{2.5pt}
  \begin{tabularx}{\linewidth}{>{\raggedright\arraybackslash}p{0.38\linewidth} >{\raggedright\arraybackslash}X}
    \toprule
    Paper label & Lean declaration path \\
    \midrule
    \path{def:main:packaged-currency} & \path{SixBirdsNeedles.Main.LegalQuotient.packagedCurrency} \\
    \path{thm:main:probe-family-currency} & \path{SixBirdsNeedles.Main.LegalQuotient.probeFamilyCurrency} \\
    \path{thm:main:minimum-spend-duality} & \path{SixBirdsNeedles.Main.LegalQuotient.minimumSpendDuality} \\
    \path{thm:main:adequacy-interface} & \path{SixBirdsNeedles.Main.XiInterface.adequacyInterface} \\
    \path{thm:main:predictive-native-membrane} & \path{SixBirdsNeedles.Main.PredictiveNative.predictiveNativeMembrane} \\
    \path{thm:main:summable-predictive-propagation} & \path{SixBirdsNeedles.Main.PredictiveNative.summablePredictivePropagation} \\
    \path{thm:main:layer-dissolving-master} & \path{SixBirdsNeedles.Main.LayerDissolving.layerDissolvingMaster} \\
    \path{def:main:involutive-ledger} & \path{SixBirdsNeedles.Main.DualityConfinement.antiInvariantObjectLedger} \\
    \path{def:main:separating-readout} & \path{SixBirdsNeedles.Main.DualityConfinement.SeparatesFixedLocus} \\
    \path{def:main:anti-invariant-membrane} & \path{SixBirdsNeedles.Main.DualityConfinement.antiInvariantCarrierCurrency} \\
    \path{lem:main:trace-identity} & \path{SixBirdsNeedles.Main.DualityConfinement.traceIdentity} \\
    \path{thm:main:separation-confinement} & \path{SixBirdsNeedles.Main.DualityConfinement.separationConfinement} \\
    \path{thm:main:quantitative-confinement} & \path{SixBirdsNeedles.Main.DualityConfinement.quantitativeConfinement} \\
    \path{def:main:domination-bridge} & \path{SixBirdsNeedles.Main.DualityConfinement.DominationBridge} \\
    \path{thm:main:douglas-domination} & \path{SixBirdsNeedles.Main.DualityConfinement.douglasDominationReverse} \\
    \path{thm:main:duality-confinement-master} & \path{SixBirdsNeedles.Main.DualityConfinement.dualityConfinementMaster} \\
    \path{def:main:exhaustive-moving-ledger} & \path{SixBirdsNeedles.Main.DualityConfinement.ExhaustiveMovingLedger} \\
    \path{thm:main:exhaustive-ledger-squeeze} & \path{SixBirdsNeedles.Main.DualityConfinement.exhaustiveLedgerSqueeze} \\
    \path{def:main:obstruction-budget} & \path{SixBirdsNeedles.Main.DualityConfinement.obstructionBudget} \\
    \path{prop:main:optimized-trace-budget} & \path{SixBirdsNeedles.Main.DualityConfinement.optimizedTraceBudget} \\
    \path{thm:main:strict-extension-contraction} & \path{SixBirdsNeedles.Main.CriticalPair.strictExtensionContraction} \\
    \path{thm:main:adequacy-aware-transfer} & \path{SixBirdsNeedles.Main.AllSix.adequacyAwareTransfer} \\
    \bottomrule
  \end{tabularx}
  \caption{Mapping from paper labels to Lean declaration paths.}
  \label{tab:paper1-lean-declaration-map}
\end{table}
 
\clearpage
\bibliographystyle{plainnat}
\bibliography{references}

\end{document}
